\documentclass[11pt]{amsart}
\usepackage{geometry}                % See geometry.pdf to learn the layout options. There are lots.
\geometry{letterpaper}                   % ... or a4paper or a5paper or ... 
%\geometry{landscape}                % Activate for for rotated page geometry
%\usepackage[parfill]{parskip}    % Activate to begin paragraphs with an empty line rather than an indent
\usepackage{graphicx}
\usepackage{amssymb}
\usepackage{epstopdf}
\DeclareGraphicsRule{.tif}{png}{.png}{`convert #1 `dirname #1`/`basename #1 .tif`.png}

\title{Problema 3.7}
%\date{}                                           % Activate to display a given date or no date

\begin{document}
\maketitle
%\section{}
%\subsection{}
\noindent  Un vendedor ha hecho una serie de ventas y desea conocer aquellas de \$200 o menos, las mayores a \$200 pero inferiores a \$400, y el número de ventas de \$400 o superiores a tal cantidad.Haga un diagrama de flujo que le proporcione al ven­dedor esta información después de haber leído los datos de entrada.  \newline
Los datos son: $N, V_1, V_2, \ldots, V_N$.   \newline

 \noindent  Donde:  \newline
 
N es una variable de tipo entero que representa el número de ventas del vendedor.
\[
\text{$V_i$}: \text{{es una variable de tipo real que indica la venta i del vendedor }} i \quad (1 \leq i \leq N)
\]
\noindent  \textbf{Explicación de las variables} \newline
I: Variable de tipo entero.Representa la variabe de control del ciclo. \newline
CHI, MED y Variables de tipo entero. Acumulan el número de ventas menores\newline
GRA:a \$200, a \$400 y mayores a \$400, respectivamente.\newline
N: Variable de tipo entero. \newline
V: Variable de tipo real. \newline

A continuación en la tabla 3.9, podemos observar el seguimiento del algoritmo.
\begin{table}[htbp]
\centering
\label{tab:corrida}
\begin{tabular}{|c|c|c|c|c|c|}
\hline
\multicolumn{6}{ |c| }{\textbf {Tabla 3.9}} \\ \hline
\textbf{l} & \textbf{M} & \textbf{MED} & \textbf{CHI} & \textbf{GRA} & \textbf{V} \\ \hline
1 & 12  & 0 &  0 &  0 & \\ \hline
2 &   & &  1&   & 180.25 \\ \hline
3 &   & &  &   1& 470.30 \\ \hline
4 &   & &  2&   & 150.25 \\ \hline
5 &   & &  3&   & 88.60 \\ \hline
6 &   & 1&  &   & 230.15 \\ \hline
7&   & &  4&   & 170.20 \\ \hline
8&   & &  5&   & 40.30 \\ \hline
9&   & 2&  &   & 201.90 \\ \hline
10&   & &  6&   & 60.32 \\ \hline
11&   &3 &  &   & 280.30 \\ \hline
12&   & &  7&   & 15.70 \\ \hline
13&   & &  8&   & 140.20 \\ \hline\end{tabular}
\end{table}

\end{document}  